\documentclass{situationist-preprint}

\title{Williams's Internalism Reconsidered}
\msid{ACT-12-b}
\msfamily{ACT-12}
\msdiscipline{Action, Agency and Free Will}
\mstype{Consequence}
\mskeywords{internal reasons; external reasons; Williams; rational criticism; bluff; explanatory constraint}

\begin{document}
\maketitle

\begin{abstract}
ACT-12-a refutes Williams's internalism: external reasons exist. We derive what survives of Williams's position. Three theses do: an agent who fails to act on an external reason commits no error of deliberation; the explanatory constraint holds in a modal form; and external reason statements cannot be vindicated to their addressee by argument from his motivations alone. None of these supports internalism, and the third shows that Williams's charge of bluff mistakes a limit on persuasion for a limit on truth.
\end{abstract}

\section{What survives}

Theorem~3.2 of ACT-12-a establishes that the Tuesday agent of \citet{Parfit1984} has a reason that is not reachable, in the sense of its Definition~2.2, and Corollary~3.3 that internalism is false. The argument relies on Williams's own account of sound deliberation (Definition~2.1 of ACT-12-a), and so leaves standing what Williams said about deliberation itself \citep{Williams1981}.

\begin{proposition}\label{p:noerror}
An agent who fails to act on an external reason, and who deliberates soundly from his motivational set, commits no error of deliberation.
\end{proposition}

\begin{proof}
By Definition~2.2 of ACT-12-a, an external reason is one to whose action no sound route leads from the agent's set. An error of deliberation is a departure from sound deliberation: a false belief retained, a means overlooked, a weighing mismanaged. The agent who deliberates soundly and does not reach the action departs from none of these.
\end{proof}

Proposition~\ref{p:noerror} locates the Tuesday agent's failure. ACT-12-a calls him irrational, and he is; but the irrationality lies in his concerns, not in his reasoning from them. The criticism that external reasons ground is criticism of what an agent cares about, and Williams was right that it cannot be presented as a correction of his deliberation.

\begin{proposition}\label{p:modal}
For every reason an agent has, it is possible for her to act for it, though not necessarily by deliberation from her present motivations.
\end{proposition}

\begin{proof}
This is the form in which Section~4 of ACT-12-a accepts the explanatory constraint: an agent may come to believe an external reason statement, and to act on it, by a route other than sound deliberation from her prior set, such as example, experience or the conversion that \citet{McDowell1995} describes. What the constraint requires is the possibility of such a route, not its deliberative character.
\end{proof}

\begin{proposition}\label{p:bluff}
An external reason statement addressed to an agent is true, but cannot be vindicated to him by argument from his motivational set alone.
\end{proposition}

\begin{proof}
It is true by Theorem~3.2 of ACT-12-a. An argument from his set that concluded in the action would be a sound route to it, and by Definition~2.2 of ACT-12-a there is none.
\end{proof}

\citet{Williams1995} suggested that external reason statements, addressed to those whom they do not move, are a kind of bluff: they claim the authority of reason while having nothing to offer the addressee's reasoning. Proposition~\ref{p:bluff} accepts the second part and rejects the first. A bluff asserts more than the speaker can back; the external reason statement asserts what is true, and what it lacks is a means of persuading the addressee by argument from his own premises. That is a limit on persuasion, not on truth.

\section{Objection and reply}

\begin{objection}
If external reasons cannot be shown to the agent by reasoning from what he cares about, they make no difference to him, and their existence is idle.
\end{objection}

\begin{reply}
They make a difference to the assessment of him, by Proposition~\ref{p:noerror} an assessment of his concerns; and, by Proposition~\ref{p:modal}, to what he may come to do. That a consideration cannot be reached by argument from one's present concerns does not make it idle any more than a truth that cannot be derived from one's present beliefs is idle.
\end{reply}

\section{Conclusion}

Internalism is false, but Williams's account of deliberation is not. An agent who ignores an external reason reasons without error; the explanatory constraint survives in a modal form; and external reason statements, though true, cannot be vindicated to their addressee from his premises. What Williams presented as a thesis about reasons is a correct thesis about deliberation and persuasion.

\begin{thebibliography}{9}\small
\bibitem[McDowell(1995)]{McDowell1995} McDowell, J. (1995). Might there be external reasons? In J.~E.~J. Altham and R. Harrison (eds.), \emph{World, Mind, and Ethics: Essays on the Ethical Philosophy of Bernard Williams}. Cambridge University Press.
\bibitem[Parfit(1984)]{Parfit1984} Parfit, D. (1984). \emph{Reasons and Persons}. Clarendon Press.
\bibitem[Williams(1981)]{Williams1981} Williams, B. (1981). Internal and external reasons. In \emph{Moral Luck: Philosophical Papers 1973--1980}. Cambridge University Press.
\bibitem[Williams(1995)]{Williams1995} Williams, B. (1995). Internal reasons and the obscurity of blame. In \emph{Making Sense of Humanity and Other Philosophical Papers 1982--1993}. Cambridge University Press.
\end{thebibliography}

\end{document}
